\documentclass[11pt,a4paper]{article}
\usepackage[T1]{fontenc}
\usepackage[utf8]{inputenc}
\usepackage[main=italian,provide=*]{babel}
\usepackage{amsmath,amssymb}
\usepackage{geometry}
\geometry{margin=2.3cm,top=2.5cm,bottom=2.7cm}
\usepackage{booktabs}
\usepackage[table]{xcolor}
\usepackage{tcolorbox}
\tcbuselibrary{skins,breakable}
\usepackage{titlesec}
\usepackage{enumitem}
\usepackage{seqsplit}
\usepackage{needspace}
\usepackage{fancyhdr}
\usepackage{hyperref}
\hypersetup{colorlinks=true,linkcolor=Sepia,citecolor=Sepia,urlcolor=Sepia}

\definecolor{inkblue}{HTML}{1B3A5C}
\definecolor{lightblue}{HTML}{EAF1F8}
\definecolor{accent}{HTML}{B0562A}
\definecolor{softgray}{HTML}{6B6B6B}
\definecolor{okgreen}{HTML}{2E6B3E}
\definecolor{okgreenbg}{HTML}{EAF5EC}
\definecolor{warnamber}{HTML}{9C6B12}
\definecolor{warnbg}{HTML}{FBF1E0}
\definecolor{advisorpurple}{HTML}{5B3A7A}
\definecolor{advisorbg}{HTML}{F1ECF6}
\definecolor{notebar}{HTML}{9C6B12}

\titleformat{\section}
  {\normalfont\Large\bfseries\color{inkblue}}
  {\thesection}{0.6em}{}[{\color{inkblue!40}\titlerule}]
\titlespacing{\section}{0pt}{0.8em}{0.4em}
\titleformat{\subsection}
  {\normalfont\large\bfseries\color{accent}}
  {}{0em}{}
\titlespacing{\subsection}{0pt}{0.5em}{0.2em}

\pagestyle{fancy}
\fancyhf{}
\renewcommand{\headrulewidth}{0.4pt}
\fancyhead[R]{\small\color{softgray}\thepage}
\fancyfoot[C]{\small\color{softgray}Tesi triennale, Samuele Schiavi --- Relatore: Prof.\ Alessandro Chiesa}

\newtcolorbox{keyeq}[1][]{
  colback=lightblue, colframe=inkblue!70, boxrule=0.6pt,
  arc=2pt, left=8pt, right=8pt, top=4pt, bottom=4pt, #1
}
\newtcolorbox{panelbox}[1]{
  colback=white, colframe=softgray!50, boxrule=0.5pt,
  arc=2pt, left=7pt, right=7pt, top=3pt, bottom=3pt,
  fonttitle=\bfseries\color{inkblue}, title={#1}
}
\newtcolorbox{okbox}[1]{
  colback=okgreenbg, colframe=okgreen!60, boxrule=0.6pt,
  arc=2pt, left=7pt, right=7pt, top=3pt, bottom=3pt,
  fonttitle=\bfseries\color{okgreen}, title={#1}
}
\newtcolorbox{warnbox}[1]{
  colback=warnbg, colframe=warnamber!70, boxrule=0.6pt,
  arc=2pt, left=7pt, right=7pt, top=3pt, bottom=3pt,
  fonttitle=\bfseries\color{warnamber}, title={#1}
}
\newtcolorbox{advisorbox}[1]{
  colback=advisorbg, colframe=advisorpurple!60, boxrule=0.6pt,
  arc=2pt, left=7pt, right=7pt, top=4pt, bottom=4pt,
  fonttitle=\bfseries\color{advisorpurple}, title={#1}
}
\newtcolorbox{editnote}{
  colback=white, colframe=white, boxrule=0pt,
  borderline west={2.2pt}{0pt}{notebar},
  arc=0pt, left=10pt, right=4pt, top=2pt, bottom=2pt,
  fontupper=\itshape\small\color{softgray!90!black}
}
\fancyhead[L]{\small\color{softgray}Shot, rumore statistico e la soglia $1/\sqrt{N_{\rm shots}}$}

\title{\vspace{-1.5em}\color{inkblue}Shot, rumore statistico\\[0.1em]e la soglia $1/\sqrt{N_{\rm shots}}$}
\author{Note di lavoro --- Tesi triennale, Samuele Schiavi\\ Relatore: Prof.\ Alessandro Chiesa}
\date{}

\begin{document}
\sloppy
\maketitle
\thispagestyle{fancy}
\vspace{-2em}

Spiegazione autosufficiente di un concetto usato più volte nei notebook
(\texttt{\seqsplit{quantum\_simulation\_dimero\_trotter.ipynb}} e
\texttt{\seqsplit{...\_semplificato.ipynb}}, sezione "Esecuzione sul
circuito"): perché si confronta l'escursione di un segnale con
$1/\sqrt{N_{\rm shots}}$, e cosa significa esattamente quel numero.
Nessuna conoscenza di statistica è data per scontata.

\section{Perché serve ripetere la misura}

Un computer quantistico, quando misuri un qubit, non ti restituisce mai
"il valore di aspettazione" --- ti restituisce \textbf{un singolo bit},
0 oppure 1, ottenuto in modo probabilistico. È una proprietà fisica
della meccanica quantistica, non un difetto dello strumento: se lo
stato è una sovrapposizione, il risultato di una singola misura è
intrinsecamente casuale, e la teoria (la "regola di Born") predice solo
\textbf{con quale probabilità} esce ciascun risultato, non quale
risultato uscirà in una singola esecuzione.

Per stimare quella probabilità --- e quindi il valore medio di un
osservabile come $\langle S_z\rangle$ --- bisogna ripetere l'intero
esperimento (preparare lo stato, farlo evolvere, misurare) molte volte
e guardare \textbf{la frazione di volte} in cui esce ciascun risultato.
Ogni ripetizione si chiama \textbf{shot}. "8192 shot" vuol dire:
l'intero esperimento è stato ripetuto 8192 volte, e si contano i
risultati.

\section{L'esempio più semplice possibile: la moneta}

Dimentica per un momento i qubit. Immagina di avere una moneta e di
volerne scoprire la probabilità di uscita testa, $p$. Se la moneta è
onesta, $p=0.5$; se è truccata, magari $p=0.503$ o $p=0.9$.

\textbf{L'unico modo per scoprirlo è lanciarla molte volte e contare.}
Se lanci la moneta $N$ volte e ottieni $k$ teste, la tua \textbf{stima}
di $p$ è
$$\hat p = \frac{k}{N}.$$
Il cappello $\hat{\ }$ sopra la $p$ è la notazione standard per dire
"stima", distinta dal vero valore $p$ che non conosci e stai cercando
di indovinare.

\textbf{Punto chiave}: $\hat p$ non è mai esattamente $p$, tranne che
per un colpo di fortuna. Ogni volta che ripeti l'intero esperimento
(altri $N$ lanci), $\hat p$ esce leggermente diverso. Questa
variabilità va sotto il nome di \textbf{rumore statistico} (o errore di
campionamento): non è un errore di misura nel senso di "lo strumento è
impreciso", è un effetto del fatto che stai cercando di dedurre una
probabilità da un numero finito di prove.

\section{Quanto è grande, di solito, questo rumore?}

Qui arriva l'unico risultato matematico di cui abbiamo bisogno, e lo si
può capire per gradi senza calcoli avanzati.

\subsection*{Intuizione: più lanci, meno rumore, ma non linearmente}

Con \textbf{pochi} lanci (es.\ $N=4$) è facilissimo ottenere
$\hat p=0.75$ (3 teste su 4) anche se la moneta è onesta ($p=0.5$): è
solo sfortuna/fortuna nel campione piccolo. Con \textbf{molti} lanci
(es.\ $N=1\,000\,000$) è quasi impossibile che $\hat p$ si discosti
molto da $0.5$: le fluttuazioni "in più" e "in meno" nei singoli lanci
tendono a compensarsi su un campione enorme.

Quindi: il rumore \textbf{diminuisce} all'aumentare di $N$. Ma di
quanto? Qui viene il punto meno intuitivo: il rumore \textbf{non}
dimezza raddoppiando $N$. Per dimezzare il rumore serve
\textbf{quadruplicare} $N$. Il rumore scala come $1/\sqrt N$, non come
$1/N$.

\subsection*{Da dove viene esattamente $1/\sqrt N$}

Ogni singolo lancio è una variabile che vale $1$ (testa) con
probabilità $p$ e $0$ (croce) con probabilità $1-p$. Un fatto di base
della probabilità (che qui prendiamo come dato, è uno dei risultati
fondazionali della teoria --- la distribuzione binomiale) dice che,
sommando $N$ lanci indipendenti di questo tipo, la variabilità del
\textbf{totale} dei successi cresce come $\sqrt N$, non come $N$.
Dividendo per $N$ per ottenere la frazione $\hat p=k/N$, la variabilità
di $\hat p$ diventa:
\begin{keyeq}
\centering
deviazione standard di $\hat p$ $=\sqrt{\dfrac{p(1-p)}{N}}$
\end{keyeq}
Questa quantità si chiama \textbf{errore standard} della stima.
"Deviazione standard" è solo un modo preciso di dire "quanto ci si
aspetta che $\hat p$ oscilli, in media, attorno al vero valore $p$,
ripetendo l'esperimento più volte".

\subsection*{Il caso peggiore: $p=0.5$}

Il fattore $p(1-p)$ nella formula dice qualcosa di interessante: è
\textbf{massimo} quando $p=0.5$ (dove vale $0.25$), e \textbf{si
azzera} quando $p\to0$ o $p\to1$.

Ha senso intuitivamente: se la moneta esce \textbf{sempre} testa
($p=1$), ogni lancio ti dà la stessa informazione --- non c'è
incertezza su cosa succederà. Se invece $p=0.5$, ogni lancio è "il più
incerto possibile" (equivale esattamente a un lancio equo), e
l'incertezza sulla stima è massima. A $p=0.5$:
$$\text{deviazione standard di }\hat p = \sqrt{\frac{0.25}{N}} = \frac{0.5}{\sqrt N}.$$

\needspace{10\baselineskip}
\section{Applicazione a un qubit e a $\langle S_z\rangle$}

Nel nostro caso non stimiamo una probabilità fra 0 e 1, ma un valore di
aspettazione $\langle S_z\rangle$ che va da $-1$ a $+1$ (autovalore
$+1$ su $|00\rangle$, $-1$ su $|11\rangle$). Con $n_{00}$ conteggi su
$N_{\rm shots}$ shot:
$$\widehat{\langle S_z\rangle} = \frac{n_{00}-n_{11}}{N_{\rm shots}} = 2\hat p - 1, \qquad \hat p \equiv \frac{n_{00}}{N_{\rm shots}}.$$
Riscalare una stima moltiplicandola per $2$ moltiplica anche la sua
deviazione standard per $2$ (la sottrazione di $1$ non cambia la
variabilità, sposta solo il centro). Quindi:
$$\text{deviazione standard di }\widehat{\langle S_z\rangle} = 2\sqrt{\frac{p(1-p)}{N_{\rm shots}}}.$$
Nel \textbf{caso peggiore} ($p=0.5$, cioè $\langle S_z\rangle=0$):
\begin{keyeq}
\centering
$2\sqrt{\dfrac{0.25}{N_{\rm shots}}} = 2\cdot\dfrac{0.5}{\sqrt{N_{\rm shots}}} = \dfrac{1}{\sqrt{N_{\rm shots}}}$
\end{keyeq}

\textbf{Ecco da dove viene esattamente il numero $1/\sqrt{N_{\rm shots}}$
usato nei notebook.} Non è un valore arbitrario né una regola empirica:
è la deviazione standard della stima di $\langle S_z\rangle$ da
campionamento, calcolata nel punto di massima incertezza possibile
($\langle S_z\rangle=0$). Con $N_{\rm shots}=8192$:
$$\frac{1}{\sqrt{8192}} \approx 0.0110.$$

\section{Perché è un caso \emph{peggiore}, e non il valore esatto ovunque}

La formula generale, $2\sqrt{p(1-p)/N_{\rm shots}}$, dipende da $p$ ---
cioè da \textbf{dove} ti trovi lungo la curva di
$\langle S_z\rangle(t)$, non solo da quanti shot usi. La tabella
seguente lo mostra con $N_{\rm shots}=8192$:

\begin{center}
\renewcommand{\arraystretch}{1.25}
\begin{tabular}{@{}cccl@{}}
\toprule
\rowcolor{inkblue!10}
$\langle S_z\rangle$ vero & $p=(\langle S_z\rangle+1)/2$ & $p(1-p)$ & deviazione standard reale\\
\midrule
$0.000$ & $0.500$ & $0.2500$ & $0.0110$ (caso peggiore, $=1/\sqrt{8192}$)\\
$0.500$ & $0.750$ & $0.1875$ & $0.0096$\\
$0.900$ & $0.950$ & $0.0900$ & $0.0066$\\
$0.994$ & $0.997$ & $0.0030$ & $0.0012$\\
$1.000$ & $1.000$ & $0.0000$ & $0.0000$ (nessuna incertezza)\\
\bottomrule
\end{tabular}
\end{center}

\begin{editnote}
Vicino agli estremi ($\langle S_z\rangle\to\pm1$) il rumore vero è
molto più piccolo del valore generico $1/\sqrt{N_{\rm shots}}$. Usare
$1/\sqrt{N_{\rm shots}}$ come soglia di riferimento è quindi una scelta
\textbf{prudenziale}: è il rumore che avresti nel punto più difficile
possibile, non necessariamente il rumore esatto nel punto che stai
guardando. È una buona regola pratica per uno screening rapido
("questo segnale è chiaramente troppo piccolo anche nel caso più
favorevole?"), non un calcolo di precisione punto per punto.
\end{editnote}

\section{Due esempi numerici concreti}

\subsection*{Un segnale che si perde nel rumore (caso generico, $p$ vicino a 0.5)}

Supponiamo di voler distinguere due valori veri di $\langle S_z\rangle$,
molto vicini fra loro: $0.030$ e $0.024$ (differenza reale: $0.006$ ---
le stesse cifre dell'escursione del set $R_0$ discusso nel notebook).
Simulando il campionamento a 8192 shot, sei ripetizioni indipendenti
dell'intero esperimento danno:

\begin{center}
\renewcommand{\arraystretch}{1.2}
\begin{tabular}{@{}cccc@{}}
\toprule
\rowcolor{inkblue!10}
ripetizione & stima (vero $=0.030$) & & stima (vero $=0.024$)\\
\midrule
1 & 0.0164 & & 0.0229\\
2 & 0.0281 & & 0.0288\\
3 & 0.0378 & & 0.0225\\
4 & 0.0112 & & 0.0210\\
5 & 0.0305 & & 0.0259\\
6 & 0.0242 & & 0.0134\\
\bottomrule
\end{tabular}
\end{center}

Le due colonne \textbf{si mescolano completamente}: guardando solo i
numeri campionati, non è possibile dire in modo affidabile quale dei
due fosse il valore vero più grande --- a volte la stima del valore
"più piccolo" (0.024) viene fuori più grande di quella del valore "più
grande" (0.030) per puro rumore (righe 2 e 3). Il rumore reale qui
($\approx0.0110$, siamo vicino a $p=0.5$) è più grande della differenza
che si vuole distinguere ($0.006$): il segnale è, di fatto, invisibile
a questo numero di shot.

\subsection*{Lo stesso segnale, ma vicino a $\langle S_z\rangle=1$ (il caso reale di $R_0$)}

Nel notebook, il set $R_0$ oscilla non attorno a $0$ ma vicino a $1$
(fra $0.994$ e $1.000$). Dalla tabella della sez.~4, lì il rumore vero
è molto più piccolo ($\sim0.0012$ invece di $0.0110$). Questo significa
che la conclusione "il segnale è invisibile" andrebbe in linea di
principio verificata punto per punto, non con la sola soglia generica.

Ciò non cambia comunque la conclusione pratica per due motivi:
\begin{enumerate}[leftmargin=1.4em,itemsep=0.25em]
\item \textbf{La soglia generica è pensata apposta per essere
conservativa}: se un segnale supera comodamente
$1/\sqrt{N_{\rm shots}}$, è sicuramente misurabile ovunque lungo la
curva; se invece è più piccolo (come $0.0062$ contro $0.0110$), serve
un controllo più attento --- che è proprio quello appena fatto nella
sez.~4, e che conferma che anche il rumore vero locale
($\sim0.001$--$0.002$ vicino agli estremi) resta dello stesso ordine di
grandezza della minuscola escursione, rendendo comunque il segnale
difficile da separare dal rumore in modo pulito su tutta la
traiettoria.
\item \textbf{Il motivo per scartare $R_0$ nel notebook non è comunque
solo questo}: la scelta del punto di lavoro richiede anche una seconda
condizione (due frequenze di oscillazione confrontabili, non solo
un'ampiezza sufficiente) --- il controllo sul rumore statistico è solo
il primo, più immediato, dei due filtri applicati.
\end{enumerate}

\needspace{6\baselineskip}
\section{Riepilogo in una frase}

\begin{okbox}{In sintesi}
$1/\sqrt{N_{\rm shots}}$ è la fluttuazione che ci si aspetta, nel caso
peggiore possibile, stimando un valore di aspettazione da
campionamento ripetuto: cresce la precisione aumentando gli shot, ma
solo con la radice quadrata, e un segnale fisico più piccolo di questa
soglia rischia di essere indistinguibile dal rumore di misura --- non
perché il calcolo o l'hardware siano imprecisi, ma per la natura
statistica della misura quantistica stessa.
\end{okbox}

\end{document}
